% Generated draft: requires mathematical and source review.
\documentclass[11pt]{article}
\usepackage[margin=25mm]{geometry}
\usepackage{fontspec}
\setmainfont{Latin Modern Roman}
\usepackage{amsmath,amssymb,mathtools,unicode-math}
\setmathfont{Latin Modern Math}
\usepackage{enumitem,needspace,xurl,hyperref}
\setlength{\parindent}{0pt}
\setlength{\parskip}{5pt}
\setlength{\emergencystretch}{3em}
\newcommand{\N}{\mathbb{N}}
\newcommand{\Z}{\mathbb{Z}}
\newcommand{\Q}{\mathbb{Q}}
\newcommand{\R}{\mathbb{R}}
\newcommand{\C}{\mathbb{C}}
\newcommand{\F}{\mathbb{F}}
\newcommand{\A}{\mathbb{A}}
\newcommand{\PP}{\mathbb{P}}
\newcommand{\E}{\mathbb{E}}
\newcommand{\eps}{\varepsilon}
\newcommand{\dd}{\,\mathrm{d}}
\begin{document}
\section*{UnsolvedMath ID: 701563}
\textbf{Source problem code:} catalogue entry
\subsection{Definitions and mathematical statement}
For a class \(M\) of groups, define
\[
L(M)=\left\{G:\text{for every }a\in G,\ \langle a^G\rangle\in M\right\},
\qquad
\langle a^G\rangle=\langle g^{-1}ag:g\in G\rangle.
\]
Here \(\langle a^G\rangle\) is the normal closure of \(a\) in \(G\); the source's phrase ``contained in \(M\)'' is understood as membership in the class \(M\).

For a finite nilpotent group \(A\), let \(qA\) be the quasivariety generated by \(A\): the class of groups satisfying every group quasi-identity valid in \(A\). A group quasi-identity is a universally quantified implication
\[
\forall x_1,\ldots,x_n\quad
\left(\bigwedge_{i=1}^{m}u_i(\mathbf{x})=v_i(\mathbf{x})\right)
\Longrightarrow u(\mathbf{x})=v(\mathbf{x}),
\]
where the terms are group terms and the conjunction is finite, possibly empty.

For a group \(H\), put \(\gamma_1(H)=H\) and \(\gamma_{k+1}(H)=[\gamma_k(H),H]\), where \([U,H]\) is generated by the commutators \(u^{-1}h^{-1}uh\) with \(u\in U\) and \(h\in H\). The group \(H\) is nilpotent if \(\gamma_{c+1}(H)=\{1\}\) for some integer \(c\geq 0\).

\textbf{Question (A. I. Budkin, Kourovka Notebook 19.15).} Does there exist a finite nilpotent group \(A\) and a group \(G\) such that
\[
\bigl(\forall a\in G,\ \langle a^G\rangle\in qA\bigr)
\quad\text{and}\quad
\bigl(\forall c\geq 0,\ \gamma_{c+1}(G)\neq\{1\}\bigr)?
\]
Equivalently, can \(L(qA)\) contain a non-nilpotent group for some finite nilpotent \(A\)? No finiteness or finite-generation assumption on \(G\) is supplied.
\subsection{Short English statement}
Can a non-nilpotent group have every normal closure of a single element in the quasivariety generated by one finite nilpotent group? The finite group may be chosen, but it must be the same for all elements.
\subsection{Sources}
\begin{itemize}
\item Local UnsolvedMath registry, permanent ID \(701563\), source code \texttt{catalogue entry}, title \emph{Kourovka Notebook Problem 19.15}; supplied upstream revision \texttt{d371ee917ac9647fbb06cd6fc529deec6f3dda82}.
\item Supplied attribution: A. I. Budkin, proposed in 2018; \emph{The Kourovka Notebook (Unsolved Problems in Group Theory)}, problem 19.15, 19th issue (2018), also identified in the supplied record with the 21st issue edition (2026), \texttt{arXiv:1401.0300v48}. The supplied collection description names E. I. Khukhro and V. D. Mazurov as editors.
\item The registry labels the problem ``open''. This is a supplied status label, not an independently verified mathematical or literature assessment. The cited editions and attribution have not been independently checked.
\end{itemize}
\end{document}
